\section{Data and Sample Construction}\label{sec:data}

\subsection{Identifying the movers}

We identify 34{,}914 workers who were employed at a convicted cartel firm during the conduct period and subsequently appeared at a different BEC-active firm after conduct ended. Of these, 2{,}273 ($6.5\%$) held managerial or professional occupations (CBO major groups 1--2); the remainder were operational staff, predominantly in service and sales roles.

The 34{,}914 movers dispersed to 3{,}960 receiving firms. The median receiving firm absorbed one ex-cartel worker; the 75th percentile absorbed three; the 90th percentile absorbed nine. We focus on the 962 receiving firms that bid on the same BEC item codes as the cartel firms both before and after absorbing ex-cartel workers, creating a balanced panel for the difference-in-differences estimation.

\subsection{Treatment and control}

\textbf{Treated firms} ($N = 962$): BEC firms that (i)~absorbed at least one worker from a convicted cartel firm after the conduct period ended, (ii)~bid on at least one item code that the cartel firm also bid on, and (iii)~have bid observations both before and after the first ex-cartel worker arrived.

\textbf{Never-treated controls} ($N = 3{,}672$): BEC firms that bid on the same item codes as the cartel firms but never absorbed any ex-cartel workers during the panel period. These firms operate in the same product markets and are exposed to the same market-level shocks, but did not receive the ``treatment'' of hiring ex-cartel workers.

\textbf{Exit-placebo controls} ($N = 3{,}335$ exit firms, $N = 10{,}183$ absorbing firms): Non-cartel BEC firms that stopped bidding between 2011 and 2014 (with at least two active years). Firms that absorbed workers from these exit firms serve as a placebo---they received a generic labor-supply shock but not a cartel-specific one. Comparing the cartel-absorber effect to the exit-absorber effect isolates the cartel-specific component of the treatment.

\subsection{Summary statistics}

Table~\ref{tab:summary} summarizes the data sources and the treatment sample.

\begin{table}[htbp]
\centering
\caption{Summary Statistics}\label{tab:summary}
\begin{tabular}{lrl}
\toprule
\multicolumn{3}{l}{\emph{Panel A: Procurement data (BEC, 2009--2019)}} \\[3pt]
Auctions (purchase-offer $\times$ item) & 837{,}706 & \\
Bids & 1{,}675{,}412 & \\
Firms & 26{,}291 & \\
Item codes & 94{,}720 & \\[6pt]
\multicolumn{3}{l}{\emph{Panel B: Employment data (RAIS, 2009--2017)}} \\[3pt]
Employment bonds & 49{,}916{,}089 & \\
Workers (unique PIS) & 15{,}283{,}903 & \\
Establishments (unique CNPJ root) & 31{,}471 & \\[6pt]
\multicolumn{3}{l}{\emph{Panel C: Cartel ground truth (CADE)}} \\[3pt]
Convicted cartels & 7 & \\
Convicted firms & 43 & (unique CNPJ roots) \\
Conduct periods & 1999--2019 & \\[6pt]
\multicolumn{3}{l}{\emph{Panel D: Worker mobility}} \\[3pt]
Workers at cartel firms during conduct & 136{,}733 & \\
Workers who moved post-conduct & 34{,}914 & \\
\quad of which managers/professionals & 2{,}273 & (6.5\%) \\
Receiving firms & 3{,}960 & \\
\quad median workers absorbed & 2 & \\
\quad 90th percentile & 15 & \\[6pt]
\multicolumn{3}{l}{\emph{Panel E: DiD estimation sample}} \\[3pt]
Treated firms (bid on cartel items pre+post) & 974 & \\
Never-treated controls & 3{,}672 & \\
Stacked DiD observations & 158{,}473 & (3 cohorts) \\
\bottomrule
\end{tabular}
\vspace{0.5em}
\begin{minipage}{0.90\textwidth}
\scriptsize{\textbf{Notes:} Panel A: preg\~{a}o auctions on BEC. Panel B: restricted to establishments whose CNPJ root appears in BEC. Panel C: SP-relevant cartels with conduct periods documented in CADE proceedings. Panel D: workers observed at a cartel firm during conduct who subsequently appear at a different BEC firm. Panel E: stacked DiD with cohorts 2012, 2014, 2015 (cohort 2013 dropped, $N = 34$).\par}
\end{minipage}
\end{table}

\subsection{Outcomes}

We study two outcomes at the firm $\times$ item-code $\times$ year level:

\begin{enumerate}
\item \textbf{Win rate}: the fraction of auctions won by the firm for a given item code in a given year. A positive treatment effect indicates that the firm became more competitive.

\item \textbf{Log price ratio}: $\ln(\text{bid price} / \text{reference price})$. A positive effect means the firm bids closer to (or above) the reference price---consistent with either market power or collusive pricing. A null or negative effect indicates competitive or unchanged pricing.
\end{enumerate}
